\chapter{Instrukcja instalacji i użytkowania}
Dostarczenie niezbędnych informacji technicznych pozwala na poprawne uruchomienie i obsługę gotowego systemu przez użytkownika końcowego. Dzięki temu proces wdrażania aplikacji w nowym środowisku staje się prosty i przewidywalny.

\section{Wymagania sprzętowe i programowe}
Aby system mógł pracować wydajnie i bez opóźnień, urządzenia wykorzystywane \\ w projekcie muszą spełniać określone parametry techniczne. Poniższe zestawienie przedstawia niezbędne minimum sprzętowe oraz oprogramowanie wymagane do poprawnej kompilacji i działania wszystkich modułów:
\begin{itemize}
	\item Komputer stacjonarny lub laptop: Procesor wspierający wirtualizację, minimum 8 GB pamięci RAM oraz system Windows 10/11, Linux lub macOS.
	\item Urządzenie mobilne: Smartfon lub tablet z systemem Android w wersji minimum 8.0.
	\item Mikrokontroler: Układ ESP32 Wroom wraz z podłączonym czujnikiem temperatury DHT11.
	\item Oprogramowanie systemowe: Docker Desktop, środowisko Java (JDK 24) oraz przeglądarka internetowa wspierająca standard HTML5.
\end{itemize}
\section{Instrukcja instalacji}
Proces instalacji systemu został podzielony na etapy, które należy wykonywać w określonej kolejności. Jest to kluczowe, ponieważ aplikacje klienckie oraz mikrokontroler wymagają aktywnego i skonfigurowanego serwera do poprawnego działania.
\begin{itemize}
	\item Uruchomienie serwera (Docker): Należy uruchomić kontener MySQL i Spring Boot za pomocą Dockera. W tym celu uruchom program Docker Desktop, otwórz folder "container things", a następnie w terminalu i wpisz komendę "docker-compose up --build".
	\item Programowanie mikrokontrolera: Korzystając z Arduino IDE, należy wgrać kod do układu ESP32. W kodzie źródłowym należy wcześniej wpisać nazwę i hasło lokalnej sieci Wi-Fi oraz adres IP komputera, na którym działa serwer.
	\item Instalacja aplikacji mobilnej: Plik aplikacji należy skompilować w Android Studio\\ i zainstalować na telefonie. W ustawieniach połączenia (klasa ApiClient/MainActivity) musi znajdować się ten sam adres IP serwera, co w przypadku ESP32.
	\item Uruchomienie interfejsu webowego: Plik index.html należy otworzyć w dowolnej nowoczesnej przeglądarce internetowej na komputerze podłączonym do tej samej sieci.
\end{itemize}

\pagebreak

\section{Instrukcja użytkowania}
Korzystanie z gotowego systemu jest intuicyjne, ponieważ większość procesów związanych z przesyłaniem i wyświetlaniem danych odbywa się automatycznie. Po prawidłowym uruchomieniu wszystkich modułów, użytkownik może monitorować temperaturę w sposób opisany poniżej.

\subsection{Obsługa aplikacji mobilnej}
Po uruchomieniu aplikacji na smartfonie, główny ekran natychmiast wyświetla ostatni zarejestrowany pomiar w wyróżnionej karcie na górze. Poniżej znajduje się chronologiczna lista wszystkich zapisów. Aplikacja sama odpytuje serwer o nowe dane co 10 sekund, więc użytkownik nie musi ręcznie odświeżać widoku, aby zobaczyć najnowsze zmiany.

\subsection{Obsługa interfejsu webowego}
W celu przejrzenia danych na komputerze wystarczy otworzyć przygotowany plik strony w dowolnej przeglądarce internetowej. Wyświetli on czytelną tabelę z kompletną historią zapisów, co pozwala na wygodną analizę danych archiwalnych.

\subsection{Kontrola stanu połączenia}
System został zaprojektowany tak, aby na bieżąco informować o ewentualnych problemach technicznych. Jeśli telefon straci połączenie z serwerem lub serwer zostanie wyłączony, na ekranie aplikacji mobilnej pojawi się komunikat informujący o braku łączności, co ułatwia szybkie zdiagnozowanie problemu z siecią Wi-Fi.